\section{External-row crossing and the precise residual in a Schur comparison}
\label{sec:r12-external-row-gaussian-residual}

\paragraph{Status and attribution.}
\textbf{PROVED THIS ROUND} are the finite unpartitioned crossing
construction, the exact residual identity (ECR9), and the Schur
matrix calculations below.  The relative primitive of the displayed
closed triple remains \textbf{OPEN}.  The sole construction
contributor is the R12 matrix agent \texttt{/root/r11\_challenge}.
This separates a genuine whole-presentation residual from an
unsupported identification with the AGC commuting-diagonal cube.

\subsection{An external row and a fixed block word}

Fix $r\ge3$, the original $K,N,\overline L,f,W$ as before, and
three selected rows, denoted $1,2,3$.  The other $r-3$ rows are
unchanged.  A block matrix $S\in\mathrm{GL}_2(\mathcal O(N))$
acts in rows $2,3$; write $B=\operatorname{diag}(1,S,I_{r-3})$.
For an analytic row vector $v=(v_2,v_3)$ fix the ordered native
word and its actual pair
\[
 V(v)=E_{12}(v_2)E_{13}(v_3),\qquad
 \mathcal P_V(v;g)
 =\mathcal P_{13,v_3}(g)
       +\mathcal P_{12,v_2}(E_{13}(v_3)g).
 \tag{ECR1}
\]
The rightmost factor acts first.  The exact matrix equality is
\[
                 BV(v)=V(vS^{-1})B.
 \tag{ECR2}
\]
Every coefficient in these words is allowed to vanish.

First supply a fixed finite unpartitioned block word
$w=B_l\cdots B_1=S$, whose factors are row units, shears in
rows $2,3$, and the fixed signed permutation $J_{23}$.
For each factor use its prescribed native pair.  There is an
explicit finite actual triple $q_w(v;g)$ with
\[
\begin{split}
 \partial_2 q_w(v;g)={}&\mathcal P_V(v;g)+p(w;V(v)g)\\
 &-p(w;g)-\mathcal P_V(vS^{-1};Bg).
\end{split}
 \tag{ECR3}
\]
This is a strict assertion about the supplied whole word, not
about a partitioned Gaussian block.

Here is a fully specified finite algorithm for the triple.  Use
the following matrix rewrite table, always retaining its displayed
native word before and after a substitution:
\[
\begin{array}{c|c|p{0.49\linewidth}}
\text{block factor}&\text{new row vector}&\text{ordered finite rules}\\\hline
D_2(\lambda)&(v_2/\lambda,v_3)
 &\text{NSC on }E_{12};\text{ passive MX on }E_{13}.\\
D_3(\lambda)&(v_2,v_3/\lambda)
 &\text{passive MX on }E_{12};\text{ NSC on }E_{13}.\\
E_{23}(c)&(v_2,v_3-cv_2)
 &\text{the signed ST substitution (ECR5) on }E_{12},E_{23};
   \text{DTC on }E_{13},E_{23};
   \text{same-target commutation and SAR to collect }E_{13}.\\
E_{32}(c)&(v_2-cv_3,v_3)
 &\text{DTC on }E_{12},E_{32};
   \text{the relabeled signed ST substitution (ECR5)};
   \text{same-target commutation and SAR to collect }E_{12}.\\
J_{23}&(-v_3,v_2)
 &\text{literal signed-permutation equivariance, followed by
   the same-target commutation of the two external shears.}
\end{array}
 \tag{ECR4}
\]
For clarity, the noncommuting shear step in the third row is the
exact consequence
\[
 E_{23}(c)E_{12}(a)
   =E_{12}(a)E_{13}(-ac)E_{23}(c)
 \tag{ECR5}
\]
of the frozen ST relation; the fourth row is its relabeling.
Its fixed substitution is the following: append the literal
inverse pair $E_{23}(-c)E_{23}(c)$ on the right, use ST with
parameters $a,-c$ on the middle $E_{12}(a)E_{23}(-c)$,
cancel the now adjacent $E_{23}(c)E_{23}(-c)$ on the left,
and commute $E_{13}(-ac)$ with $E_{12}(a)$ by the specified
same-target witness.  The resulting word is (ECR5).
The collection uses the fixed order of (ECR1).  For every rule
take precisely its frozen full triple, with its orientation
reversed when the rule is used backwards, evaluated at the
chronological starting tuple of that subword.  Common prefixes
and suffixes are retained exactly as in GB's finite substitution
rule.  Each coefficient in (ECR4) is an analytic polynomial in
the old coefficients and the supplied unit inverses.

The same-target commutation in this table has a precise current
source interface: expand its four native shears into the eight
unit edges and use the SAN triple $\mathcal J_0$ for that closed
loop.  Every edge changes only row $1$, all other entries remain
literally the same, the norm vertices telescope and the unit
weights sum to zero.  SAN therefore applies with normalization
zero.  Its prescribed isometry jumps have zero native pair.
Every individual shear companion is zero on its own
small-coefficient open by TCT2, so the complete expression
extends through all coefficient zeros.  This is a derived
same-target facet; it does not restore the missing historical
CG/RA source bytes or count that historical relation as a new
independent cycle-existence class.

To make the starting tuples in the algorithm unambiguous, put
$g_0=g$, $v^{(0)}=v$, and recursively
\[
        g_j=B_jg_{j-1},\qquad
        v^{(j)}=v^{(j-1)}B_j^{-1}.
 \tag{ECR6}
\]
Apply the appropriate row of (ECR4) at $(g_{j-1},v^{(j-1)})$
and add all its oriented actual triples.  Then sum over
$j=1,\ldots,l$.  The intermediate external path pairs cancel
literally, giving exactly (ECR3).  All rules in the table are
strict.  In particular the passive MX rule uses its full
difference of CS witnesses with canceling extending symbols,
not either CS witness alone.  No cap is changed: each call keeps
its own entire frozen expression and zero extension.  Thus this
is an explicit finite construction with every metric and mixed
coefficient term, on the original ordinary domains.

\subsection{What changes after the GB partition}

Now take the two exact GB words $w_\alpha,w_\gamma$ for $S$,
their original overlap triple $Q_S(g)$, and their original
ordinary partition $\theta$.  All these coefficient data depend
on $S$, not on the input tuple $g$.  The extending pair in GB18
is also independent of $g$.  Use (ECR3) to construct the two
whole triples $q_\alpha(v;g),q_\gamma(v;g)$ without changing
any of their frozen caps.  On the punctured pivot overlap set
\[
 \mathcal T(S,v;g)=q_\alpha(v;g)-q_\gamma(v;g)
                        -Q_S(V(v)g)+Q_S(g).
 \tag{ECR7}
\]
This is a specified finite actual triple.  It includes all of
the ST, DTC, passive-MX, same-target and GB cap terms just listed.
The two copies of $Q_S$ are evaluated at their \emph{different}
original tuples, so their passive subtuples have not been
identified.  Subtracting the two exact identities (ECR3) and
then using GB18 proves
\[
                       \partial_2\mathcal T(S,v;g)=0.
 \tag{ECR8}
\]
The GB errors cancel because their coefficients and original
metric are identical, not because one cancels a curvature factor.

Let $p_S(g)$ be the entire partitioned GB pair, and define the
external crossing pair
\[
 \mathcal X(S,v;g)=\mathcal P_V(v;g)+p_S(V(v)g)
                      -p_S(g)-\mathcal P_V(vS^{-1};Bg).
\]
The complete partition calculation is
\[
 \boxed{\quad
 \mathcal X(S,v;g)
 =\partial_2\bigl(\theta q_\alpha+(1-\theta)q_\gamma\bigr)
                      -K_2(\theta)\mathcal T(S,v;g).
 \quad}
 \tag{ECR9}
\]
Indeed $p_S=\theta p_\alpha+(1-\theta)p_\gamma+K_2(\theta)Q_S$.
Insert this identity twice and use the full product differential.
Its derivative terms are precisely
$-K_2(\theta)(q_\alpha-q_\gamma)
 +K_2(\theta)(Q_S(V(v)g)-Q_S(g))$, giving (ECR9).
All products have the closed-support zero extensions specified
by RAG.  No term containing $d'\theta$ or $d''\theta$ has
been omitted.  The pair $\mathcal X$ is independently
$\partial_1$-closed, as also follows from GB's input-independent
curvature expression and (ECR2).

For reference, the \emph{exact sufficient missing statement} is
a relative identity on an ordinary neighborhood of the closed
supports of the partition derivatives,
\[
              \mathcal T=\partial_3 H+e_2|_W,
         \qquad H\text{ actual},\quad e_2\in L_2^{r-1}(N)
                       \text{ actual on that neighborhood}.
 \tag{ECR10}
\]
If it is supplied, the full derivative-insertion identity
$\partial_2K_3+K_2\partial_3=0$ yields the explicit certificate
\[
 \mathcal X
   =\partial_2\bigl(\theta q_\alpha+(1-\theta)q_\gamma
                                  +K_3(\theta)H\bigr)
                               -K_2(\theta)e_2|_W.
 \tag{ECR11}
\]
The final error extends as an actual pair on $N$ by closed-support
zero extension.  In this section (ECR10) is \textbf{OPEN}, not a
consequence of (ECR8).  In particular one cannot normalize every
ST and GB witness to a cap based only on the last $r-3$ rows
using TCT: their original unchanged passive blocks include
different members among rows $1,2,3$.  TCT explicitly supplies
no such cross-passive-block assertion.  Nor has (ECR7) been
identified with the six-face raw cube of AGC.

\subsection{Why this residual occurs already at three Gaussian pivots}

Take an arbitrary analytic invertible three-row block
\[
 M=\begin{pmatrix}a&b&c\\d&e&f\\g&h&i\end{pmatrix}
\]
inside an arbitrary tuple of length $r\ge3$.  On $a\ne0$ put
$t=d/a$, $u=g/a$, and
\[
 S_a=\begin{pmatrix}E&F\\H&I\end{pmatrix}
 =\begin{pmatrix}e-tb&f-tc\\h-ub&i-uc\end{pmatrix},\qquad
 U_a=E_{12}(b/a)E_{13}(c/a).
 \tag{ECR12}
\]
The first Gaussian factorization is
\[
       M=E_{31}(u)E_{21}(t)
               D_1(a)\operatorname{diag}(1,S_a)U_a.
 \tag{ECR13}
\]
For the first pivots $d$ and $g$, use the signed permutations
$J_{12}$ and $J_{13}$, respectively, then the same rule.  On
the corresponding overlaps direct matrix multiplication gives
\[
 \begin{split}
 S_d&=\begin{pmatrix}t^{-1}&0\\-u/t&1\end{pmatrix}S_a
                         =E_{32}(-u)D_2(t^{-1})S_a,\\
 S_g&=\begin{pmatrix}1&-t/u\\0&u^{-1}\end{pmatrix}S_a
                         =E_{23}(-t)D_3(u^{-1})S_a.
 \end{split}
 \tag{ECR14}
\]
The row labels in the right formulas are the ambient labels
$2,3$.  On $adg\ne0$ their transition satisfies
\[
 T_{ga}T_{da}^{-1}
  =\begin{pmatrix}0&-t/u\\1&1/u\end{pmatrix}
                  =J_{23}D_3(t/u)E_{23}(u^{-1}).
 \tag{ECR15}
\]
These are concrete two-row triangular/unit/Weyl words, rather
than an abstract invocation of Gaussian elimination.

They are not the entire comparison of the native pairs.  The
upper prefixes also change:
\[
 \begin{split}
 U_d&=V\bigl((E,F)/d\bigr)U_a,
       &\operatorname{diag}(1,S_a)U_d
          &=E_{12}(d^{-1})\operatorname{diag}(1,S_a)U_a,\\
 U_g&=V\bigl((H,I)/g\bigr)U_a,
       &\operatorname{diag}(1,S_a)U_g
          &=E_{13}(g^{-1})\operatorname{diag}(1,S_a)U_a.
 \end{split}
 \tag{ECR16}
\]
Thus the specific residuals (ECR7) occur with $S=S_a$ and
$v=(E,F)/d$ or $(H,I)/g$, starting at the actual tuple $U_af$.
The row-unit and triangular multiplication theorems address the
coefficient-matrix part (ECR14); they do not delete (ECR16).
There may in addition be higher associativity requirements for
chosen pair-comparison triples.  No complete three-pivot triple
compatibility theorem is inferred from these matrix identities.

\paragraph{Exact boundary of this calculation.}
This source provides both a constructive strict crossing for every
supplied finite unpartitioned block word and the full, explicitly
located obstruction after GB partitioning.  The obstruction is
not asserted nonzero, nonextendable, or a counterexample to Strong G.
It is also not replaced by selected Gauss chambers, tensor probes
or an abstract quotient primitive.  A further construction of
(ECR10), with all the displayed original caps and ordinary zero
domains, would be needed before this route gives a whole
$\mathrm{GL}_3$ block certificate or an original-$U$ actual
Green form.  The same external-row interface recurs in the
$\mathrm{GL}_k$ Schur recursion; low-dimensional notation here
describes that general mechanism, not a lower-dimensional
restriction of the final Green problem.

\paragraph{Exact frozen inputs.}
SAN: \path{c8416ab1d14618f8f23724bcdad19db08a36928cb19b590a46d2306701a7a3b6}.
TCT: \path{effdf9332b6fa606237a99a25217cdfdd44b7fd09cd685f654ff221201d68c2c}.
ST: \path{da4de041961b99f38c6378f448a52cf30b875a122384bbcb3f207e745be8a7ca}.
DTC: \path{9bf95b4b9ae7ec0717ec573b135c390172efcadfd5ff2b00242c2a0e7f648117}.
MX: \path{ba5eebf4b312b7df4f3a7beceacaafbfaa171cbe97eb2614af7e717197d450fa}.
GB: \path{f4f87d6aabfc325f7d2f01d629a6f4a9fa80227cb9600cf592ea380b0b5104d9}.
NSC, SAR, RAG and SHT are the exact inputs named in these frozen
sources.  The AGC source is used only to distinguish its narrower
fixed-cube statement from the unresolved identity (ECR10).
